% SGA 3，第 VII 篇_B，英文重建本——第 1 轮（中文译本）。
% 范围：依据原文件第 8--11 页，包含 0.3 至 0.3.4；止于原文件第 11 页的 0.3.5 之前。
% 合并阅读版页码： 521--524. 再版印刷页码： 510--513.
% 本范围内的原版页边定位： 482--484.
% 依据： Polo--Gille Exp7B-13oct24.pdf,
% SHA-256 D303690E0195654BBE7C3DAA5D89EA23E24CD709A3575217234AC19AE8836694.
% 曾参阅 Jacob Reinhold 的 CC BY 4.0 Markdown，作为对比版本来源。
% 本范围内无原生图；所有陈列数学均为原生 LaTeX。

\medskip
\noindent\textbf{0.3.}\footnote{编者注：本段已经修改，以计入 \SGARefLink{sga3:VIIB:0.2}{0.2} 中所作的补充。}
\phantomsection
\label{sga3:VIIB:0.3}
设 \(L\) 和 \(M\) 是两个伪紧 \(A\)-模。函子
\[
  \operatorname{LF}(A)\longrightarrow(\mathrm{Ab}),
  \qquad
  N\longmapsto\operatorname{Bil}_{c}(L\times M,N),
\]
是左正合的，其中 \(\operatorname{Bil}_{c}(L\times M,N)\) 表示从
\(L\times M\) 到 \(N\) 的连续 \(A\)-双线性映射组成的阿贝尔群。

因此，由 \SGARefLink{sga3:VIIB:0.2.E}{0.2.E}，存在表示此函子的伪紧 \(A\)-模
\(L\widehat{\otimes}_{A}M\)，它在唯一同构意义下唯一；即对于
\(\operatorname{LF}(A)\) 的每个对象 \(N\)，存在函子性同构
\[
  \operatorname{Hom}_{c}
    \bigl(L\widehat{\otimes}_{A}M,N\bigr)
  \simeq
  \operatorname{Bil}_{c}(L\times M,N).
\]

此外，\(L\widehat{\otimes}_{A}M\) 等同于投射极限
\[
  P=\widehat{L\otimes_{A}M}
\]
其中取极限的对象是离散 \(A\)-模
\[
  (L/L')\otimes_{A}(M/M'),
\]
其中 \(L'\) 和 \(M'\) 分别遍历 \(L\) 和 \(M\) 的开子模。

事实上，设 \(\varphi:L\times M\to N\) 是到有限长度的离散
\(A\)-模 \(N\) 的连续双线性映射。由下文的
\SGARefLink{sga3:VIIB:0.3.1}{引理~0.3.1}，存在 \(L\) 和 \(M\) 的开子模
\(L'\) 和 \(M'\)，使得
\[
  \varphi(L'\times M)=\varphi(L\times M')=\{0\}.
\]
这意味着由 \(\varphi\) 诱导的映射
\[
  \overline{\varphi}:L\otimes_{A}M\longrightarrow N
\]
具有形式 \(\varphi'\circ p\)，其中 \(p\) 是从 \(L\otimes_{A}M\) 到
\((L/L')\otimes_{A}(M/M')\) 的典范投影。若以 \(\widehat{\varphi}\) 表示复合映射
\[
  P
  \longrightarrow
  (L/L')\otimes_{A}(M/M')
  \xrightarrow{\ \varphi'\ }
  N,
\]
则映射 \(\varphi\mapsto\widehat{\varphi}\) 是从
\(\operatorname{Bil}_{c}(L\times M,N)\) 到 \(\operatorname{Hom}_{c}(P,N)\)
的函子性双射。因此
\[
  P\simeq L\widehat{\otimes}_{A}M.
\]

因此，伪紧模 \(L\widehat{\otimes}_{A}M\) 是 \(L\otimes_{A}M\)
对于由典范投影的核所定义的线性拓扑的分离完备化，这些投影为
\[
  L\otimes_{A}M
  \longrightarrow
  (L/L')\otimes_{A}(M/M'),
\]
我们称它为 \(L\) 与 \(M\) 的\emph{完备化张量积}。

若 \(x\in L\) 且 \(y\in M\)，则将 \(x\otimes_{A}y\) 在
\(L\widehat{\otimes}_{A}M\) 中的像记为 \(x\widehat{\otimes}_{A}y\)。

\medskip
\noindent\textbf{引理 0.3.1。}
\phantomsection
\label{sga3:VIIB:0.3.1}
\textit{设 \(L\)、\(M\) 和 \(N\) 是伪紧 \(A\)-模，其中 \(N\) 的长度有限。
若 \(\varphi:L\times M\to N\) 是连续的 \(A\)-双线性映射，则存在
\(L\) 和 \(M\) 的开子模 \(L'\) 和 \(M'\)，使得}
\[
  \varphi(L'\times M)=\varphi(L\times M')=\{0\}.
\]

\noindent\textit{证明。}
逆像 \(\varphi^{-1}(0)\) 是 \((0,0)\) 的开邻域，故包含形如
\(L_{1}\times M_{1}\) 的开集，其中 \(L_{1}\) 和 \(M_{1}\) 是
\(L\) 和 \(M\) 的开子模。由于 \(L/L_{1}\) 长度有限，
存在 \(L\) 的元素 \(x_{1},\ldots,x_{r}\)，使得
\[
  L_{1}+Ax_{1}+\cdots+Ax_{r}=L.
\]
若 \(M'\subset M_{1}\) 充分小，则对每个 \(i\) 都有
\(\varphi(x_i,M')=0\)，因为映射 \(y\mapsto\varphi(x_i,y)\) 连续。
由此得到 \(\varphi(L,M')=\{0\}\)。同样，
当 \(L'\) 充分小时，\(\varphi(L',M)=\{0\}\)。
\hfill\(\square\)

\medskip
\noindent\textbf{推论 0.3.1.1。}\footnote{编者注：本推论为新增。}
\phantomsection
\label{sga3:VIIB:0.3.1.1}
\textit{设 \(M\) 是伪紧 \(A\)-模。}
\begin{enumerate}
  \item[\textup{(i)}]
  \textit{对每个开子模 \(M'\)，存在 \(A\) 的开理想 \(\mathfrak l\)，使得}
  \[
    \overline{\mathfrak lM}\subset M'.
  \]

  \item[\textup{(ii)}]
  \textit{因此，}
  \[
    M\simeq
    \varprojlim_{\mathfrak l}
      M/\overline{\mathfrak lM},
  \]
  \textit{其中 \(\mathfrak l\) 遍历 \(A\) 的开理想所构成的滤过投射系，
  每个 \(M/\overline{\mathfrak lM}\) 都赋予使其成为伪紧模的商拓扑；
  参见\ \SGARefLink{sga3:VIIB:0.2.D}{0.2.D}。}
\end{enumerate}

\noindent\textit{证明。}
考虑映射
\[
  \varphi:A\times M\longrightarrow M/M',
  \qquad
  (a,m)\longmapsto am+M'.
\]
由 \SGARefLink{sga3:VIIB:0.3.1}{0.3.1}，存在 \(A\) 的开理想
\(\mathfrak l\)，使 \(\mathfrak lM\subset M'\)。由于 \(M'\) 也是闭的，
它包含 \(\overline{\mathfrak lM}\)。由于 \(M\) 的所有开子模的交为零，得到
\[
  \bigcap_{\mathfrak l}\overline{\mathfrak lM}=(0).
\]
另一方面，由 \SGARefLink{sga3:VIIB:0.2.D}{0.2.D}，映射
\[
  \phi:M\longrightarrow
  \varprojlim_{\mathfrak l}M/\overline{\mathfrak lM}
\]
是满射。因此，由 \SGARefLink{sga3:vii:proposition:0-2-b}{0.2.B} 的证明，
\(\phi\) 诱导同构
\[
  M/\operatorname{Ker}(\phi)
  \xrightarrow{\ \sim\ }
  \varprojlim_{\mathfrak l}M/\overline{\mathfrak lM}.
\]
而我们刚刚已经看到
\[
  \operatorname{Ker}(\phi)
  =
  \bigcap_{\mathfrak l}\overline{\mathfrak lM}
  =
  (0).
\]
\hfill\(\square\)

\medskip
\noindent\textbf{注记 0.3.1.2。}\footnote{编者注：本注记为新增。}
\phantomsection
\label{sga3:VIIB:0.3.1.2}
完备化张量积满足通常的结合律：若 \(L\)、\(M\) 和 \(P\) 是伪紧
\(A\)-模，则有典范同构
\[
  (L\widehat{\otimes}_{A}M)\widehat{\otimes}_{A}P
  \simeq
  L\widehat{\otimes}_{A}(M\widehat{\otimes}_{A}P).
\]
事实上，两个对象都表示将 \(\operatorname{LF}(A)\) 的每个对象 \(N\)
对应到从 \(L\times M\times P\) 到 \(N\) 的连续 \(A\)-三线性映射
所成的阿贝尔群的函子。

\medskip
\noindent\textbf{0.3.2.}
\phantomsection
\label{sga3:VIIB:0.3.2}
设
\[
  L'\xrightarrow{\,f\,}L\xrightarrow{\,g\,}L''\longrightarrow0
\]
是正合列，且 \(M\) 是 \(\operatorname{PC}(A)\) 的对象。显然，对于
\(\operatorname{LF}(A)\) 的每个对象 \(N\)，所诱导的序列
\[
\begin{array}{ccccccc}
0 &\longrightarrow&
  \operatorname{Bil}_{c}(L''\times M,N)
&\longrightarrow&
  \operatorname{Bil}_{c}(L\times M,N)
&\longrightarrow&
  \operatorname{Bil}_{c}(L'\times M,N)
\\[0.35em]
&&\Big\Vert&&\Big\Vert&&\Big\Vert
\\[-0.1em]
0 &\longrightarrow&
  \operatorname{Hom}_{c}(L''\widehat{\otimes}_{A}M,N)
&\longrightarrow&
  \operatorname{Hom}_{c}(L\widehat{\otimes}_{A}M,N)
&\longrightarrow&
  \operatorname{Hom}_{c}(L'\widehat{\otimes}_{A}M,N)
\end{array}
\]
都是正合的。由 \SGARefLink{sga3:VIIB:0.2.E}{0.2.E}，这等价于说序列
\begin{equation}
  L'\widehat{\otimes}_{A}M
  \xrightarrow{\,f\widehat{\otimes}M\,}
  L\widehat{\otimes}_{A}M
  \xrightarrow{\,g\widehat{\otimes}M\,}
  L''\widehat{\otimes}_{A}M
  \longrightarrow0
  \tag{\(*\)}
  \label{sga3:VIIB:0.3.2.star}
\end{equation}
正合。因此：

\medskip
\noindent\textbf{推论。}
\phantomsection
\label{sga3:VIIB:0.3.2.cor}
\textit{对于每个伪紧 \(A\)-模 \(M\)，函子 \(-\widehat{\otimes}_{A}M\) 右正合。}

特别地，取 \(L=A\)，令 \(f\) 是闭理想 \(\mathfrak a\) 到 \(A\) 的包含映射，
令 \(g\) 是 \(A\) 到 \(A/\mathfrak a\) 的典范投影。于是可以通过映射
\[
  x\widehat{\otimes}_{A}m\longmapsto xm.
\]
将 \(A\widehat{\otimes}_{A}M\) 与 \(M\) 等同。
由于 \(\mathfrak a\widehat{\otimes}_{A}M\) 的像在 \(M\) 中闭
（参见\ \SGARefLink{sga3:vii:proposition:0-2-b}{0.2.B}），且
\(\mathfrak a\otimes_{A}M\) 的像在 \(\mathfrak a\widehat{\otimes}_{A}M\) 中稠密，
故 \(f\widehat{\otimes}M\) 的像恰是 \(\mathfrak aM\) 在 \(M\) 中的闭包
\(\overline{\mathfrak aM}\)。因此，正合列 \(\textup{(\(*\))}\) 给出同构
\[
  (A/\mathfrak a)\widehat{\otimes}_{A}M
  \xrightarrow{\ \sim\ }
  M/\overline{\mathfrak aM}.
\]

\medskip
\noindent\textbf{0.3.3. Nakayama 引理。}
\phantomsection
\label{sga3:VIIB:0.3.3}
\textit{设 \(A\) 是伪紧环，\(M\) 是伪紧 \(A\)-模，\(\mathfrak a\) 是 \(A\)
中包含于根 \(r(A)\) 的理想。则}
\[
  \overline{\mathfrak aM}=M
  \qquad\Longrightarrow\qquad
  M=0.
\]

\noindent\textit{证明。}
假设
\(\overline{\mathfrak aM}=M\).\footnote{编者注：必要时将 \(\mathfrak a\) 换成其闭包，可以假定 \(\mathfrak a\) 是闭的。}
设 \(M'\) 是 \(M\) 的开子模，并令 \(M''=M/M'\)。
由于 \(M''\) 是离散的，\(\mathfrak aM''\) 在 \(M''\) 中闭，
故等于 \(\overline{\mathfrak aM''}\)。由 \SGARefLink{sga3:VIIB:0.3.2}{0.3.2}，典范映射
\[
  M/\overline{\mathfrak aM}
  \longrightarrow
  M''/\overline{\mathfrak aM''}
\]
是满射，所以
\[
  \mathfrak aM''
  =
  \overline{\mathfrak aM''}
  =
  M''.
\]
由于 \(M''\) 是有限生成 \(A\)-模，且 \(\mathfrak a\subset r(A)\)，
通常的 Nakayama 引理给出 \(M''=0\)。因此，\(M\) 的每个开子模
\(M'\) 都等于 \(M\)，所以 \(M=0\)。\footnote{编者注：事实上，\(M\) 是各商 \(M/M'=0\) 的投射极限。}
\hfill\(\square\)

\medskip
\noindent\textbf{0.3.4.}
\phantomsection
\label{sga3:VIIB:0.3.4}
由此得到 Nakayama 引理的通常推论。

\medskip
\noindent\textbf{推论。}
\phantomsection
\label{sga3:VIIB:0.3.4.cor}
\textit{设 \(\mathfrak a\) 是包含于 \(r(A)\) 的闭理想，且
\(f:M\to N\) 是伪紧 \(A\)-模的态射。}
\begin{enumerate}
  \item[\textup{(i)}]
  \textit{若诱导映射}
  \[
    f':
    M/\overline{\mathfrak aM}
    \longrightarrow
    N/\overline{\mathfrak aN}
  \]
  \textit{是满射，则映射 \(f\) 是满射。}\footnote{编者注：因此，每个伪紧 \(A\)-模都是一个拓扑自由 \(A\)-模的商。
  首先归约到 \(A\) 为局部环的情形；参见\ \SGARefLink{sga3:VIIB:0.3.7}{0.3.7} 的证明。}

  \item[\textup{(ii)}]
  \textit{若 \(N\) 是投射对象，则只要 \(f'\) 是同构，\(f\) 就是同构。}
\end{enumerate}

\noindent\textit{证明。}
将\SGARefLink{sga3:VIIB:0.3.3}{引理~0.3.3} 应用于
\(\operatorname{Coker}(f)\)，即得 \textup{\SGARefLink{sga3:vii:paragraph:5-1:equation:i}{(i)}}。
对于 \textup{\SGARefLink{sga3:vii:paragraph:5-1:equation:ii}{(ii)}}，设 \(f'\) 是同构。
由 \textup{\SGARefLink{sga3:vii:paragraph:5-1:equation:i}{(i)}}，\(f\) 是满射，从而有截面。
再将\SGARefLink{sga3:VIIB:0.3.3}{引理~0.3.3} 应用于 \(\operatorname{Ker}(f)\)。
\hfill\(\square\)\par\smallskip

当 \(A\) 是极大理想为 \(\mathfrak m\) 的局部环时，还可由
\SGARefLink{sga3:VIIB:0.3.3}{0.3.3} 推出下列交换定理。

\medskip
\noindent\textbf{定理。}
\phantomsection
\label{sga3:VIIB:0.3.4.exchange}
\textit{设 \(A\) 是局部伪紧环，\(\mathfrak m\) 是其极大理想，\(M\) 是具有伪基
\((m_i)_{i\in I}\) 的拓扑自由 \(A\)-模
\textup{（\SGARefLink{sga3:VIIB:0.2.1}{0.2.1}）}，且 \(N\) 是 \(M\) 的闭直和项。
则 \(M\) 有一个由 \(N\) 中的元素及部分 \(m_i\) 组成的伪基。}

\noindent\textit{证明。}
当 \(A\) 是域时，这一点显然：使用 \SGARefLink{sga3:VIIB:0.2.2}{0.2.2} 的对偶性，
并应用通常的交换定理。因此，
\[
  N/\overline{\mathfrak mN}
\]
有一个拓扑自由的补 \(A/\mathfrak m\)-模，具有伪基
\((\overline{m}_i)_{i\in J}\)，其中 \(\overline{m}_i\) 是 \(m_i\) 在
\[
  M/\overline{\mathfrak mM},
\]
中的像，且 \(J\) 是 \(I\) 的子集。\footnote{编者注：我们已将 \(N/\mathfrak mN\) 改为 \(N/\overline{\mathfrak mN}\)，
下文的 \(M/\overline{\mathfrak mM}\) 也作了相应修改。}
若
\[
  P=\prod_{i\in J}Am_i,
\]
则典范映射
\[
  N\oplus P\longrightarrow M
\]
“模 \(\mathfrak m\) 是双射”。因此，由上文，它是双射。
另一证明见 \textup{[CA]}，\S~IV.2，命题~8。
\hfill\(\square\)

% 下一断点： Exp7B 原文件页码 11, 合并阅读版页码 524,
% 再版印刷页码 513, 紧接其前： 0.3.5.
